
\documentclass{article}
\usepackage{colortbl}
\usepackage{makecell}
\usepackage{multirow}
\usepackage{supertabular}

\begin{document}

\newcounter{utterance}

\centering \large Interaction Transcript for game `dond', experiment `semi\_en', episode 18 with Qwen3{-}32B{-}t0.0.
\vspace{24pt}

{ \footnotesize  \setcounter{utterance}{1}
\setlength{\tabcolsep}{0pt}
\begin{supertabular}{c@{$\;$}|p{.15\linewidth}@{}p{.15\linewidth}p{.15\linewidth}p{.15\linewidth}p{.15\linewidth}p{.15\linewidth}}
    \# & $\;$A & \multicolumn{4}{c}{Game Master} & $\;\:$B\\
    \hline

    \theutterance \stepcounter{utterance}  
    & & \multicolumn{4}{p{0.6\linewidth}}{
        \cellcolor[rgb]{0.9,0.9,0.9}{
            \makecell[{{p{\linewidth}}}]{
                \texttt{\tiny{[P1$\langle$GM]}}
                \texttt{You are playing a negotiation game in which you have to agree on how to divide a set of items among you and another player.} \\
\\ 
\texttt{Rules:} \\
\texttt{(a) You and the other player are given a set of items. Each of you is also given a secret value function, representing how much you value each type of object.} \\
\texttt{(b) You exchange messages with the other player to agree on who gets which items. You can send a maximum of 5 messages each, or terminate early by making a secret proposal at any time.} \\
\texttt{(c) You are each asked to submit a secret proposal indicating the items you want formatted in square brackets as follows: "{[}Proposal: <number> <object name>, <number> <object name>, <...>{]}"} \\
\texttt{(d) If your proposals are complementary, i.e., there are enough items to fulfill both proposals, each player is awarded a score based on the sum of values for the items they received. Otherwise, both of you get zero points.} \\
\texttt{(e) Your goal is to maximize the score you receive.} \\
\\ 
\texttt{Let us start.} \\
\\ 
\texttt{The set of available items is:} \\
\\ 
\texttt{1 stick, 1 towel, 1 toy, 2 boxes} \\
\\ 
\texttt{Your secret value function is:} \\
\\ 
\texttt{stick: 5, towel: 2, toy: 1, box: 1} \\
\\ 
\texttt{IMPORTANT: Your messages, unless it is the secret proposal, are directly transmitted to the other player, so do not include any response to the rules or text announcing your message. To make a secret proposal, use the indicated format. Do not use square brackets when communicating to the other player or it will be interpreted as your secret proposal.} \\
            }
        }
    }
    & \\ \\

    \theutterance \stepcounter{utterance}  
    & \multicolumn{4}{p{0.6\linewidth}}{
        \cellcolor[rgb]{0.9,0.9,0.9}{
            \makecell[{{p{\linewidth}}}]{
                \texttt{\tiny{[P1$\rangle$GM]}}
                \texttt{<think>} \\
\texttt{Okay, let's see. I need to negotiate with another player to divide the items: 1 stick, 1 towel, 1 toy, and 2 boxes. My value function is stick:5, towel:2, toy:1, box:1. The other player has their own secret values, which I don't know. The goal is to maximize my score by making sure our proposals are complementary.} \\
\\ 
\texttt{First, I should figure out what the other player might value. Since I don't know their values, I need to make some assumptions. Maybe they value the stick highly too, or maybe they prefer other items. But since I have to maximize my own points, I should prioritize getting the stick because it's worth 5 for me. The next is the towel at 2, then the toy and boxes at 1 each.} \\
\\ 
\texttt{But if I take the stick, maybe the other player can take the towel and the boxes. Wait, there are two boxes. If I take the stick, and the other takes the towel and both boxes, that would be complementary. Let me check: my proposal would be {[}Proposal: 1 stick{]}, and their proposal would be {[}Proposal: 1 towel, 2 box{]}. That way, all items are allocated without overlap. My score would be 5, and theirs would be 2 + 2*1 = 4. Since the proposals are complementary, we both get our scores. That seems good for me.} \\
\\ 
\texttt{Alternatively, if I take the stick and a box, then the other player could take the towel and the remaining box. My score would be 5 +1=6, and theirs 2+1=3. But maybe the other player would prefer to take more boxes? But since I don't know their values, maybe they might want the stick. If I take the stick, they might not want to take nothing. But if I take the stick and a box, maybe they can take the towel and the other box. That's still complementary. But I need to make sure that my proposal doesn't conflict with theirs.} \\
\\ 
\texttt{Wait, the problem is that if I take the stick, and they also propose to take the stick, then the proposals are not complementary, and we both get zero. So I need to make sure that my proposal doesn't overlap with theirs. But how can I know what they want? Maybe I should propose taking the stick and then see if they can} \\
            }
        }
    }
    & & \\ \\

    \theutterance \stepcounter{utterance}  
    & & & \multicolumn{2}{p{0.3\linewidth}}{
        \cellcolor[rgb]{0.9,0.9,0.9}{
            \makecell[{{p{\linewidth}}}]{
                \texttt{\tiny{[GM$|$GM]}}
                \texttt{proposal} \\
            }
        }
    }
    & & \\ \\

    \theutterance \stepcounter{utterance}  
    & & & \multicolumn{2}{p{0.3\linewidth}}{
        \cellcolor[rgb]{0.9,0.9,0.9}{
            \makecell[{{p{\linewidth}}}]{
                \texttt{\tiny{[GM$|$GM]}}
                [1, 0, 0, 0]
            }
        }
    }
    & & \\ \\

    \theutterance \stepcounter{utterance}  
    & & \multicolumn{4}{p{0.6\linewidth}}{
        \cellcolor[rgb]{0.9,0.9,0.9}{
            \makecell[{{p{\linewidth}}}]{
                \texttt{\tiny{[GM$\rangle$P2]}}
                \texttt{You are playing a negotiation game in which you have to agree on how to divide a set of items among you and another player.} \\
\\ 
\texttt{Rules:} \\
\texttt{(a) You and the other player are given a set of items. Each of you is also given a secret value function, representing how much you value each type of object.} \\
\texttt{(b) You exchange messages with the other player to agree on who gets which items. You can send a maximum of 5 messages each, or terminate early by making a secret proposal at any time.} \\
\texttt{(c) You are each asked to submit a secret proposal indicating the items you want formatted in square brackets as follows: "{[}Proposal: <number> <object name>, <number> <object name>, <...>{]}"} \\
\texttt{(d) If your proposals are complementary, i.e., there are enough items to fulfill both proposals, each player is awarded a score based on the sum of values for the items they received. Otherwise, both of you get zero points.} \\
\texttt{(e) Your goal is to maximize the score you receive.} \\
\\ 
\texttt{Let us start.} \\
\\ 
\texttt{The set of available items is:} \\
\\ 
\texttt{1 stick, 1 towel, 1 toy, 2 boxes} \\
\\ 
\texttt{Your secret value function is:} \\
\\ 
\texttt{stick: 7, towel: 3, toy: 0, box: 0} \\
\\ 
\texttt{IMPORTANT: Your messages, unless it is the secret proposal, are directly transmitted to the other player, so do not include any response to the rules or text announcing your message. To make a secret proposal, use the indicated format. Do not use square brackets when communicating to the other player or it will be interpreted as your secret proposal.} \\
\\ 
\\ 
\texttt{The other player has submitted a secret proposal. It is now your turn to submit a proposal of your own.} \\
            }
        }
    }
    & \\ \\

    \theutterance \stepcounter{utterance}  
    & & & \multicolumn{4}{p{0.6\linewidth}}{
        \cellcolor[rgb]{0.9,0.9,0.9}{
            \makecell[{{p{\linewidth}}}]{
                \texttt{\tiny{[GM$\langle$P2]}}
                \texttt{<think>} \\
\texttt{Okay, let's see. I need to figure out what proposal to make here. The available items are 1 stick, 1 towel, 1 toy, and 2 boxes. My value function is stick:7, towel:3, toy:0, box:0. So I care most about the stick, then the towel. The other player has already submitted a secret proposal, so I need to respond with my own.} \\
\\ 
\texttt{First, I should consider what the other player might have proposed. Since their proposal is secret, I don't know exactly, but I can make some educated guesses. Let's think about possible scenarios. If they asked for the stick, then I need to make sure that my proposal doesn't conflict with that. But since the stick is the highest value for me, I probably want it. However, if the other player also wants the stick, then our proposals would conflict, and we both get zero. So maybe I should check if they want the stick. But since they've already made a proposal, I can't ask them directly. I have to make my proposal based on possible combinations.} \\
\\ 
\texttt{Alternatively, maybe the other player's proposal is for items that don't include the stick. For example, if they want the towel or the toy or boxes. Since I don't value the toy or boxes, maybe they do. Let's think about possible complementary proposals. If I take the stick and the towel, that's two items. The other player could take the toy and both boxes. But there are two boxes, so maybe they want multiple boxes. But since there are two boxes, if they want two boxes, that's possible. Let me check the numbers.} \\
\\ 
\texttt{The total items are 1 stick, 1 towel, 1 toy, 2 boxes. If I propose to take the stick and the towel, that's 2 items. The other player could take the toy and both boxes, which is 3 items. But there are only 5 items total? Wait, no. Wait, the items are 1 stick, 1 towel, 1 toy, 2 boxes. That's 1+1+1+2=5 items. If I take stick and towel (2 items), and the other takes toy and 2 boxes (3 items), that adds up to 5. So that would be complementary. So if I take stick and towel, and the other takes the rest, then our} \\
            }
        }
    }
     \\ \\

    \theutterance \stepcounter{utterance}  
    & & & \multicolumn{2}{p{0.3\linewidth}}{
        \cellcolor[rgb]{0.9,0.9,0.9}{
            \makecell[{{p{\linewidth}}}]{
                \texttt{\tiny{[GM$|$GM]}}
                \texttt{continue} \\
            }
        }
    }
    & & \\ \\

    \theutterance \stepcounter{utterance}  
    & & & \multicolumn{2}{p{0.3\linewidth}}{
        \cellcolor[rgb]{0.9,0.9,0.9}{
            \makecell[{{p{\linewidth}}}]{
                \texttt{\tiny{[GM$|$GM]}}
                \texttt{player was instructed to make a proposal} \\
            }
        }
    }
    & & \\ \\

\end{supertabular}
}

\end{document}
